% sec-03: independence and simultaneity.  Table 8, Figure 5 (mermaid-type gantt).
\section{Independence and Simultaneity}

Two words in the request carry most of its weight. \emph{Independently} means
the company had no relationship with Revolution Medicines, no access to its data,
and no right of reference to its filings; every input to the company's work was
public. \emph{Simultaneously} means the work was done while the developer's
trials were running, not after their results were known. Both are claims about
time, and both can be checked against public registries. Table 8 sets the two
calendars side by side.

\begin{apptable}
\begin{tabularx}{\textwidth}{>{\raggedright\arraybackslash}p{2.9cm} >{\raggedright\arraybackslash}p{6.3cm} >{\raggedright\arraybackslash}X}
\headrow \thead{Period} & \thead{Revolution Medicines, public record} & \thead{ChemicalQDevice, deposited record}\\
\midrule
2022 to 2024 & Phase 1/1b opens; RASolute 302 opens & No work on the molecule \\
June 2025 & RASolute 302 enrolling & Identification across forty meta-analyses \\
August 2025 & RASolute 302 enrolling & QSP simulation, 12.8 against 5.4 months \\
December 2025 & RASolute 304, adjuvant Phase 3, opens & Physical AI work follows in 2026 \\
May to July 2026 & Readout: 13.2 against 6.6 months & Protocols, IND and funding applications \\
August 26, 2026 & FDA approval as Rasonque & Ten applications, capitalization plan \\
September 2026 & RASolute 304 enrolling & Site package and these funding actions \\
\end{tabularx}
\end{apptable}
\vspace{-0.60cm}
\tabcap{Table 8. The developer's trial calendar and the company's deposit calendar on one\\
set of rows, showing work begun while the Phase 3 trial was still enrolling patients.}

The developer's rows come from the public registries and the published readout
\cite{ctgov001,ctgov302,ctgov304,rasolute302}. The company's rows come from its
own deposit dates \cite{daraxident,qspmetpancre,onpremwhippl,phase1ind,movein2026}.
Nothing in either column depends on the other. Figure 5 draws the same two
calendars as a gantt chart, so the overlap can be seen rather than read.

\begin{appfloat}[!tb]
\begin{appfig}[mermaid-type / gantt chart, two sections on one axis]
% Axis: mid 2024 to late 2026 at 5 cm per year, x = (year - 2024.5) * 5.
% Developer bars are gray; company bars and milestones are in the accent.
\fill[fundgrayl] (2.50,0.40) rectangle (7.50,-4.25);
\foreach \x/\lab in {0/{mid 2024}, 2.50/2025, 7.50/2026, 11.25/{Q4 2026}}{
  \draw[fundgray,line width=0.3pt] (\x,0.40) -- (\x,-4.25);
  \node[font=\tiny\sffamily,text=fundgray,anchor=south] at (\x,0.42) {\lab};}
\node[mmlanetitle,anchor=west] at (-3.30,0.05) {Revolution Medicines};
\ganttrow{-0.40}{0.00}{11.50}{mmbarg}{Phase 1/1b, since 2022}
\ganttrow{-0.95}{1.50}{9.50}{mmbarg}{RASolute 302, Phase 3}
\ganttrow{-1.50}{7.25}{11.50}{mmbarg}{RASolute 304, adjuvant}
\node[anchor=east,font=\tiny\sffamily] at (-0.12,-2.05) {Readout and approval};
\node[diamond,draw=fundink,fill=fundgrayd,minimum size=2.6mm,inner sep=0pt] at (9.55,-2.05) {};
\node[font=\tiny\sffamily,anchor=east] at (9.35,-2.05) {302 reported};
\node[diamond,draw=fundink,fill=fundaccent,minimum size=2.6mm,inner sep=0pt] at (10.75,-2.05) {};
\node[font=\tiny\sffamily,anchor=west] at (10.95,-2.05) {Approved};
\node[mmlanetitle,anchor=west] at (-3.30,-2.60) {ChemicalQDevice};
\node[anchor=east,font=\tiny\sffamily] at (-0.12,-3.05) {2025 papers};
\node[diamond,draw=fundink,fill=fundaccent,minimum size=2.6mm,inner sep=0pt] at (4.75,-3.05) {};
\node[diamond,draw=fundink,fill=fundaccent,minimum size=2.6mm,inner sep=0pt] at (5.60,-3.05) {};
\node[font=\tiny\sffamily,anchor=west] at (5.85,-3.05) {Identification, then QSP simulation};
\ganttrow{-3.60}{9.80}{10.30}{mmbark}{Protocols, IND, applications}
\ganttrow{-4.15}{10.40}{11.25}{mmbark}{Plans, site package, funding blocks}
\draw[fundaccent,dashed,line width=0.7pt] (10.75,0.40) -- (10.75,-4.25);
\pnote{-3.30}{-4.70}{Shaded band: calendar 2025. Dashed line: the approval on August 26, 2026.}
\end{appfig}
\vspace{-0.60cm}
\figcaption{Figure 5. The developer's three trials and the company's deposits on one time\\
axis, with the company's first two papers inside the Phase 3 enrollment window.}
\end{appfloat}

\subsection{What the overlap shows, and what it does not}

The overlap shows that the company chose the molecule and simulated its effect
while the Phase 3 trial was enrolling, and published both before the trial
reported. It does not show that the simulation predicted the trial. The
simulation modeled 1000 patients on a combination in a KRAS G12C population; the
trial enrolled 241 patients in the relevant comparison on a single agent in a
primarily G12D and G12V population. The two survival figures, 12.8 and 13.2
months, are close, and the company describes that closeness as a dated chronology
and never as a validation \cite{qspmetpancre,rasolute302}.

It also does not show that the developer knew of the company's work or relied on
it. There is no evidence either way, and the company claims none. Independence
runs in both directions: the company did not use the developer's data, and it
does not suggest the developer used its papers.
